\documentclass[10pt,a4paper]{amsart}
\usepackage[margin=19mm]{geometry}
\usepackage{amsmath,amssymb,mathtools}
\usepackage[scr=rsfs,cal=euler]{mathalfa}
\usepackage{xfrac}
\usepackage[unicode,hidelinks,bookmarks=false]{hyperref}
\newcommand{\ie}{i.e.\ }
\newcommand{\cf}{cf.\ }
\newcommand{\resp}{respectively\ }
\newcommand{\Subsection}{\subsection}
\providecommand{\nobreakdash}{\nobreak\hbox{-}}
\setlength{\emergencystretch}{2em}
\multlinegap0pt
\usepackage{amssymb}
\usepackage{enumerate}
\usepackage{dsfont}
\newtheorem{prop}{Proposition}
\DeclareMathOperator{\Leb}{Leb}
\title{A Ruelle-McMullen formula for the volume dimension\\ of skew products in $\mathbb C^2$}
\author{Fabrizio Bianchi, Yan Mary He}
\date{Source: arXiv:2603.06475v1 (CC BY 4.0)}
\begin{document}
\maketitle

\noindent\emph{Source and licence.} A Ruelle-McMullen formula for the volume dimension\\ of skew products in $\mathbb C^2$. Version arXiv:2603.06475v1 (CC BY 4.0), distributed by arXiv under the Creative Commons Attribution 4.0 International licence, http://creativecommons.org/licenses/by/4.0/, as recorded in the arXiv licence metadata for this exact version and shown on the abstract page https://arxiv.org/abs/2603.06475v1. Full licence notice: \texttt{licenses/arxiv-2603.06475-CC-BY-4.0.txt}.\par\smallskip

\noindent\textit{Editorial scope.} Counted unit: the article's prop prop\_extra\_term (source0.tex lines 1027--1040). The article's complete original proof of it is delivered with it (source0.tex lines 1086--1190, one \texttt{proof} environment of 105 lines, which the article places away from the statement). No mathematics is written, repaired or re-derived: the statement and the proof are the article's own lines, in the article's own notation, and no retained line is substituted or re-flowed.  Retained uncounted as the source-local context the proof needs: lines 1051--1057.  Nothing was omitted from any retained span, so the package carries no omission record.  No retained line contains a citation, so the wrapper carries no bibliography and no bibliography entry.  No retained line contains a figure environment, an image-inclusion command or drawing code, so no image is delivered and the package declares no asset dependency.  Editorial formatting only, in addition: an AMS \texttt{amsart} preamble that restates the article's own declarations and macros used by the retained lines, namely the environments \texttt{prop} on separate counters, together with the standard packages those lines need.  The retained environments print in this document as Proposition 1 in order of appearance, whereas the article prints the same items with the numbering of its own full document.  The complete native source of the article as distributed in the arXiv e-print for this exact version is bundled unchanged as \texttt{sources/195877/source0.tex}.
\begin{equation}\label{eq_v_outside}
v_z(w)
=
-\frac{w}{d}
\sum_{k=1}^d\sum_{n=0}^{\infty}
\frac{c_k(z^{d^n})}{d^n}\,w^{-k d^n}.
\end{equation}

\begin{prop}\label{prop_extra_term}
We have
\[
I
\left(
\frac{\partial v}{\partial w}\right)
=
\frac{1}{d^2\log d
}
\int_{S^1}
\sum_{k=1}^d k^2\,|c_k(z)|^2\,
d\Leb_1 (z).
\]
\end{prop}

\begin{proof}[Proof of Proposition \ref{prop_extra_term}]
Differentiating~\eqref{eq_v_outside}
 term by term gives
\[
\frac{\partial v}{\partial w}(z,w)
=
\frac{1}{d}
\sum_{k=1}^d\sum_{n=0}^{\infty}
\left(
\frac{c_k(z^{d^n})}{d^n}\,w^{-k d^n}
+
k\,c_k(z^{d^n})\,w^{-k d^n-1}
\right).
\]
The terms involving the denominator $d^{n}
$ are uniformly summable and negligible in the
logarithmic limit defining $I(\partial_w v)$.
Hence only the second sum contributes, and we obtain
\[
I\!\left(\frac{\partial v}{\partial w}\right)
=
\lim_{r\to1^+}
\frac{1}{|\log(r-1)|}
\int_{S^1}
\int_{|w|=r}
\left|
\frac{1}{d}
\sum_{k=1}^d
\sum_{n=0}^{\infty}
k\,c_k(z^{d^n})\,w^{-k d^n-1}
\right|^2
\, 
d\Leb_r(w)\,
d\Leb_1(z).
\]

Since all exponents $-k d^n-1$ are distinct, the functions
$w\mapsto w^{-k d^n-1}$
 are orthogonal in $L^2(\Leb_r)$.
Therefore, for every $z\in S^1$, we have
\[
\int_{|w|=r}
\left|
\frac{1}{d}
\sum_{k=1}^d
\sum_{n=0}^{\infty}
k\,c_k(
z^{d^n}
)\,w^{-k d^n-1}
\right|^2
\,
d\Leb_r (w)
=
\frac{1}{d^2}
\sum_{k=1}^d\sum_{n=0}^{\infty}
k^2|c_k(z^{d^n})|^2\,|r|^{-2k d^n-2}.
\]
For any fixed $r>1$,
 the factor $|r|^{-2k d^n-2}$
  is close to $1$ for
$n\le N(r)$ and exponentially small for $n>N(r)$, where
\[
N(r)\sim\frac{|\log(r-1)|}{\log d}.
\]
Hence,
\[
\int_{|w|=r}
\left|
\frac{1}{d}
\sum_{k=1}^d
\sum_{n=0}^{\infty}
k\,c_k(z^{d^n})\,w^{-k d^n-1}
\right|^2
\,
d\Leb_r (w)
=
\frac{1}{d^2}
\sum_{k=1}^d\sum_{n=0}^{N(r)}
k^2|c_k(z^{d^n})|^2
+
o(|\log(r-1)|).
\]

Dividing by $|\log(r-1)|$ and letting $r\to1$, we may apply the Birkhoff
ergodic theorem to the sequence $\{z^{d^n}\}_{n\in \mathbb N}$.
For 
$\Leb_1$-almost every $z\in S^1$
 we obtain
\[
\lim_{r\to1^+}
\frac{1}{|\log(r-1)|}
\int_{|w|=r}
\left|\frac{\partial v}{\partial w}(z,w)\right|^2\,
d\Leb_r (w)
=
\frac{1}{d^2\log d}
\sum_{k=1}^d
k^2\int_{S^1}|c_k(z')|^2\,
d\Leb_1(z')
\]

Integrating with respect to $z$
(observe that the right hand side is constant)
 yields the assertion.
\end{proof}

\end{document}
